\section{生物学示例族 II：工作记忆与吸引子维持}

把吸引子语言引入本手稿的最清晰理由之一，来自理论神经科学中关于工作记忆的标准讨论。
持续活动与类吸引子维持模型之所以有用，
是因为它们解释了：即便触发线索已经消失，系统如何仍能足够久地保留方向性，从而维持任务。

\begin{definition}[类积分器持续性]
\textbf{类积分器}机制，是任何能够保存任务相关信息足够长时间、从而偏置后续状态空间运动的结构。
\end{definition}

\begin{discussion}
这一点会以实践方式影响本书的框架。
一个计划、一个笔记，或一个外部提醒，都可以被看作粗糙的类积分器辅助。
它保存了本来会在转变完成之前衰减掉的方向性信息。
数学并不能证明一张写下来的便签在字面上就是一个皮层积分器，
但它确实给了读者一个比“更努力去记住”更严肃的概念。
\end{discussion}

\begin{example}[写作场景中的维持]
一名研究生打开草稿，在失去势头之前写下三条 bullet claims。
如果这三条主张仍然保持可见，它们就可以充当一种类积分器式记忆痕迹，
使系统继续停留在写作盆地附近。
如果它们消失不见，状态就可能重新漂移回阅读或浏览。
因此，工作记忆类比并不是行为证明；它是对“为什么可见的持续性辅助重要”的一种中等车道数学澄清。
\end{example}